\chapter{Pendahuluan}
\label{bab:pendahuluan}

\section{Latar Belakang}

Deskripsikan latar belakang dari permasalahan yang akan diangkat pada
penelitian tesis. Latar belakang berisi penjelasan dari \textit{Problem Domain}
yang termuat pada judul penelitian. Misalkan penulis mengambil suatu judul
berikut "Deteksi Penyakit Kanker Dengan Sistem Pakar Berbasis Fuzzy". Judul
tersebut mempunyai Problem adalah Deteksi Penyakit Kanker, Problem Domain adalah
Penyakit Kanker, dan Uniqueness adalah Sistem Pakar Berbasis Fuzzy.

Berkaitan dengan Latar Belakang, Problem Domain-nya adalah tentang penyakit
kanker, sehingga disini penulis dapat menceritakan tentang penyakit kanker dan
tingkat urgensi (seperti tingkat kenaikan jumlah penderitanya dari tahun ke
tahun, gawatnya penyakit kanker tersebut, dan semisalnya). Latar belakang yang
baik berisikan problem domain yang mempunyai tingkat urgensi tinggi.

\section{Permasalahan}

Deskripsikan dengan jelas permasalahan yang ingin diteliti pada tesis.
Permasalahan berisi penjelasan dari \textit{Problem} yang termuat pada judul
penelitian. Deskripsi masalah sebaiknya dituliskan dengan gaya bahasa
deskriptif. Pada contoh judul diatas, Problem-nya adalah tentang deteksi
penyakit kanker, sehingga penulis disini dapat menceritakan tentang sulitnya
pendeteksian penyakit kanker dengan mendeskripsikan faktor-faktor yang membuat
sulit dalam pendeteksiannya. Uraian permasalahan yang baik manakala penulis
berhasil menyakinkan kepada pembaca tentang seberapa tingkat urgensi dari
permasalahan tersebut sehingga membuat pembaca yakin bahwa permasalahan
tersebut membutuhkan solusi dari penelitian pada tesis penulis.

\section{Tujuan}

Deskripsikan dengan jelas tujuan penelitian tesis yang diangkat. Tujuan tesis
harus jelas, singkat, dan mengandung klaim orisinalitas. Tuliskan secara
argumentatif apa saja fitur-fitur yang ditawarkan pada penelitian sebagai
sesuatu solusi yang baru untuk mengklaim orisinalitas pada penelitian tesis.
Untuk memberikan gambaran yang jelas apa yang dilakukan dan solusi unik apa
yang akan ditawarkan pada penelitian, tujuan sebaiknya diawali dengan kalimat
pembuka seperti ini: ``\textit{Penelitian tesis ini mengajukan suatu pendekatan
/ algoritma / pemodelan / teknik / metode yang baru untuk mengatasi
........................} (hubungan dengan Problem) \textit{dengan memakai /
menggunakan / mempresentasikan / melibatkan .....................} (sebutkan
solusi memakai apa)''. Kalimat-kalimat selanjutnya kemudian memperjelas solusi
dengan fitur-fitur unik seperti apa yang ditawarkan pada penelitian sebagai
suatu solusi untuk menjawab permasalahan.

Pada contoh judul diatas, orisinalitasnya adalah Sistem Pakar Berbasis Fuzzy,
sehingga penulis disini dapat menjelaskan tentang pemodelan sistem pakar dan
fitur-fitur fuzzy yang seperti apa untuk deteksi penyakit kanker. Kalimat
pertama pada tujuan dapat diawali dengan contoh berikut, ``\textit{Penelitian
tesis ini mengajukan suatu pendekatan baru untuk pendeteksian penyakit kanker
dengan mempresentasikan model klasifikasi menggunakan Sistem Pakar yang
berbasis Fuzzy}''. Kemudian pada kalimat-kalimat berikutnya terangkan secara
argumentatif tentang fitur-fitur unik pemodelan Sistem Pakar dengan Fuzzy
sehingga dapat digunakan untuk mendeteksi penyakit kanker.

\section{Manfaat}

Uraikan kontribusi tesis pada pengembangan ilmu pengetahuan teknologi dan seni,
pemecahan masalah pembangunan, atau pengembangan kelembagaan. Kontribusi
menggambarkan manfaat dari penelitian terhadap pihak tertentu saat penelitian
sudah selesai. Kontribusi sebaiknya bersifat spesifik, tidak terlalu luas dan
tidak terkesan mengada-ada. Jelaskan siapa yang mendapatkan manfaat dari
penelitian penulis dan dalam bentuk apa manfaatnya.

\section{Sistematika Penulisan}

Jelaskan tentang sistematika pembahasan dalam buku tesis yang meliputi:

\begin{sistematika}
  \babsistematika{Bab 1}{Pendahuluan}{Jelaskan tentang apa saja yang dibahas
    pada Bab 1. Penjelasan memuat bagian-bagian penting pada Pendahuluan.}
  \babsistematika{Bab 2}{Kajian Pustaka}{Jelaskan tentang apa saja yang dibahas
    pada Bab 2. Penjelasan memuat bagian-bagian penting pada Kajian Pustaka.}
  \babsistematika{Bab 3}{Desain Sistem}{Jelaskan tentang apa saja yang dibahas
    pada Bab 3. Penjelasan memuat bagian-bagian penting pada Desain Sistem.}
  \babsistematika{Bab 4}{Eksperimen dan Analisis}{Jelaskan tentang apa saja
    yang dibahas pada Bab 4. Penjelasan memuat bagian-bagian penting pada
    Eksperimen dan Analisis.}
  \babsistematika{Bab 5}{Penutup}{Jelaskan tentang apa saja yang dibahas pada
    Bab 5. Penjelasan memuat bagian-bagian penting pada Penutup.}
\end{sistematika}
